\documentclass[11pt,a4paper]{article}
\usepackage[margin=2.5cm,headheight=15pt]{geometry}
\usepackage{amsmath,amssymb,amsthm}
\usepackage{mathtools}
\usepackage{microtype}
\usepackage{hyperref}
\usepackage{cleveref}
\usepackage{enumitem}
\usepackage{mdframed}
\usepackage{xcolor}
\usepackage{titlesec}
\usepackage{booktabs}
\usepackage{fancyhdr}

\definecolor{navyblue}{RGB}{26,58,92}
\definecolor{midblue}{RGB}{44,95,138}
\definecolor{lightblue}{RGB}{232,240,248}
\definecolor{proofgray}{RGB}{242,242,242}
\definecolor{opencolor}{RGB}{235,245,235}
\definecolor{openborder}{RGB}{0,120,0}
\definecolor{propcolor}{RGB}{240,248,240}
\definecolor{propborder}{RGB}{0,130,0}

\hypersetup{colorlinks=true,linkcolor=midblue,citecolor=midblue,urlcolor=midblue}

\pagestyle{fancy}\fancyhf{}
\fancyhead[L]{\small\color{navyblue}\textit{SSP Variants that Collapse to JGP}}
\fancyhead[R]{\small\color{navyblue}\thepage}
\renewcommand{\headrulewidth}{0.4pt}

\titleformat{\section}{\large\bfseries\color{navyblue}}{}{0em}{}[\color{navyblue}\titlerule]
\titleformat{\subsection}{\normalsize\bfseries\color{midblue}}{}{0em}{}
\titleformat{\subsubsection}{\normalsize\itshape\color{midblue}}{}{0em}{}

\newmdenv[backgroundcolor=lightblue,linecolor=navyblue,linewidth=1.2pt,
  leftline=true,rightline=false,topline=false,bottomline=false,
  innerleftmargin=12pt,innerrightmargin=8pt,innertopmargin=8pt,
  innerbottommargin=8pt,skipabove=10pt,skipbelow=6pt]{thmbox}

\newmdenv[backgroundcolor=proofgray,linecolor=midblue,linewidth=0.8pt,
  leftline=true,rightline=false,topline=false,bottomline=false,
  innerleftmargin=10pt,innerrightmargin=8pt,innertopmargin=6pt,
  innerbottommargin=6pt,skipabove=6pt,skipbelow=6pt]{proofbox}

\newmdenv[backgroundcolor=opencolor,linecolor=openborder,linewidth=1.2pt,
  leftline=true,rightline=false,topline=false,bottomline=false,
  innerleftmargin=12pt,innerrightmargin=8pt,innertopmargin=8pt,
  innerbottommargin=8pt,skipabove=10pt,skipbelow=6pt]{openproblem}

\newmdenv[backgroundcolor=propcolor,linecolor=propborder,linewidth=1.2pt,
  leftline=true,rightline=false,topline=false,bottomline=false,
  innerleftmargin=12pt,innerrightmargin=8pt,innertopmargin=8pt,
  innerbottommargin=8pt,skipabove=10pt,skipbelow=6pt]{propbox}

\theoremstyle{plain}
\newtheorem{theorem}{Theorem}[section]
\newtheorem{lemma}[theorem]{Lemma}
\newtheorem{proposition}[theorem]{Proposition}
\newtheorem{corollary}[theorem]{Corollary}
\theoremstyle{definition}
\newtheorem{definition}[theorem]{Definition}
\newtheorem{remark}[theorem]{Remark}

\DeclareMathOperator{\OPT}{OPT}
\newcommand{\Kstar}{K^*}
\newcommand{\ZJGP}{Z_{\mathrm{JGP}}^*}
\newcommand{\ZSSP}{Z_{\mathrm{SSP}}^*}
\newcommand{\calG}{\mathcal{G}}

\setlist[itemize]{noitemsep,topsep=4pt}
\setlist[enumerate]{noitemsep,topsep=4pt}

% ─────────────────────────────────────────────────────────────────────────────
\begin{document}
% ─────────────────────────────────────────────────────────────────────────────

\begin{center}
  {\color{navyblue}\rule{\textwidth}{2pt}}\\[0.4cm]
  {\LARGE\bfseries\color{navyblue}The Job Sequencing and Tool Switching Problem}\\[0.3cm]
  {\large\itshape\color{midblue}Part VII: SSP Variants that Collapse to JGP}\\[0.2cm]
  {\color{navyblue}\rule{\textwidth}{2pt}}
\end{center}

\begin{mdframed}[backgroundcolor=lightblue,linecolor=navyblue,linewidth=1pt,
  innerleftmargin=14pt,innerrightmargin=14pt,innertopmargin=8pt,innerbottommargin=8pt]
\small\textbf{Scope and Status.}
This document addresses the question: can we design a variant of the SSP with
a modified cost function $f$ such that solving the variant \emph{collapses} to
solving the JGP? We present:
(i)~a formal definition of collapse;
(ii)~the TSIP as a known baseline case (\textbf{proved in literature});
(iii)~\textbf{sufficient conditions} for collapse (C1+C2+C3) — \textbf{proved};
(iv)~the magazine setup time variant $S_\mathrm{stop}$ and when it causes
collapse — \textbf{proved};
(v)~a \textbf{necessity analysis} of the conditions (\cref{sec:analysis}): (C3) is
necessary, while (C1) and (C2) are not, each with an explicit witness.
The document explicitly distinguishes proved from open results.
\end{mdframed}

\vspace{0.5cm}
\tableofcontents
\newpage

% ─────────────────────────────────────────────────────────────────────────────
\section{What Does ``Collapse to JGP'' Mean?}
\label{sec:collapse_def}
% ─────────────────────────────────────────────────────────────────────────────

\begin{definition}[SSP-$f$ variant]
Given a cost function $f: \mathcal{C} \times \mathcal{C} \to \mathbb{R}_{\ge 0}$,
the SSP-$f$ variant minimises $\sum_{k} f(C_k, C_{k+1})$ over all feasible
job sequences and magazine configurations.
\end{definition}

\begin{definition}[Collapse to JGP]
\label{def:collapse}
SSP-$f$ \emph{collapses to JGP} if for every instance $I$:
\[
  \OPT_{\mathrm{SSP}\text{-}f}(I)
  \;\text{ is achieved by a solution using exactly } \ZJGP(I) \text{ groups.}
\]
Equivalently: using more than $\ZJGP$ groups never strictly reduces the
SSP-$f$ objective.
\end{definition}

\noindent If collapse holds, the JGP can be solved first and the SSP-$f$ objective
minimised within those groups, dramatically simplifying the problem.

% ─────────────────────────────────────────────────────────────────────────────
\section{Baseline: The Tool Switching Instants Problem (TSIP)}
\label{sec:tsip}
% ─────────────────────────────────────────────────────────────────────────────

\begin{definition}[TSIP cost function]
$f_{01}(C_k, C_{k+1}) = \mathbf{1}[C_k \ne C_{k+1}]$: the cost is 1 if the
configuration changes, 0 otherwise.
\end{definition}

The TSIP minimises the total number of \emph{configuration changes} (regardless
of how many tools are swapped each time). With $K$ groups, there are at most
$K-1$ configuration changes, so the TSIP objective is $\le K-1$.

\begin{thmbox}
\begin{theorem}[TSIP collapses to JGP \cite{Adjiashvili2015}]
\label{thm:tsip}
Under $f_{01}$, the optimal SSP-$f_{01}$ solution uses exactly $\ZJGP$ groups.
\end{theorem}
\end{thmbox}

\begin{proofbox}
\begin{proof}
Any solution with $K$ configuration changes corresponds to processing $K$
groups (one per distinct configuration used). So $\OPT_{\mathrm{TSIP}} = \ZJGP - 1$
(minimum configuration changes = minimum groups minus one). \qed
\end{proof}
\end{proofbox}

\begin{remark}
The TSIP is solvable in polynomial time when tool loading/unloading is unconstrained
\cite{Adjiashvili2015}. It models situations where the machine must stop completely
for each tool change (the ``batch'' model).
\end{remark}

% ─────────────────────────────────────────────────────────────────────────────
\section{Sufficient Conditions for Collapse (Proved)}
\label{sec:sufficient}
% ─────────────────────────────────────────────────────────────────────────────

We give general conditions under which a cost function $f$ causes SSP-$f$ to
collapse to JGP.

\begin{thmbox}
\begin{theorem}[Sufficient conditions for collapse]
\label{thm:sufficient}
SSP-$f$ collapses to JGP if $f$ satisfies all three of the following conditions:
\begin{description}
\item[(C1) Self-cost zero:] $f(C, C) = 0$ for all $C \in \mathcal{C}$.
\item[(C2) Strict inter-group positivity:] There exists $\delta > 0$ such that
  for any JGP-optimal grouping $\calG^*$, any two groups $G_k \ne G_{k'}$, and
  any valid configurations $C \in \mathcal{C}(G_k)$, $C' \in \mathcal{C}(G_{k'})$:
  $f(C, C') \ge \delta > 0$.
\item[(C3) No-beneficial-split:] For every feasible grouping $\calG$ with
  $K > \Kstar$ groups, the SSP-$f$ cost for any solution over $\calG$ is $\ge$
  the SSP-$f$ cost for the best $\Kstar$-group solution.
\end{description}
\end{theorem}
\end{thmbox}

\begin{proofbox}
\begin{proof}
Under (C2), any solution with $K$ groups incurs $\ge K-1$ inter-group transitions,
each costing $\ge \delta > 0$, so total cost $\ge (K-1)\delta$.
For $K > \Kstar$: $(K-1)\delta > (\Kstar-1)\delta$, meaning a $\Kstar$-group
solution with cost $(\Kstar-1)\delta$ (e.g., by using the TSIP baseline) is
strictly cheaper than any $(K+1)$-group solution.
Condition (C3) ensures that no splitting of groups into more pieces can further
reduce cost. Together, every optimal solution under SSP-$f$ uses exactly $\Kstar$
groups. \qed
\end{proof}
\end{proofbox}

\begin{remark}[Necessity — partial answer, with an open core]
The conjunction (C1)+(C2)+(C3) is \emph{sufficient} but not necessary: the necessity
analysis of \cref{sec:analysis} (\cref{prop:necessity}) shows (C3) is necessary while
(C1) and (C2) are each dispensable, via explicit witnesses (e.g.\ collapse can hold with
positive self-cost, violating (C1), or with some zero-cost inter-group transitions,
violating (C2)). What remains open is a single \emph{necessary and sufficient}
characterisation of collapse; \cref{sec:analysis} reduces this to the threshold
$\rho_c(I)$ but does not close it.
\end{remark}

% ─────────────────────────────────────────────────────────────────────────────
\section{The Magazine Setup Time Variant (Proved)}
\label{sec:setup_time}
% ─────────────────────────────────────────────────────────────────────────────

\subsection{Model}

In many FMS environments, changing any tool requires stopping the machine.
This stop incurs a fixed overhead $S_\mathrm{stop}$ (machine re-initialisation,
operator time) in addition to the per-tool switching cost $c_\mathrm{switch}$.
The total cost of a transition from $C_k$ to $C_{k+1}$ is:
\[
  f_{S}(C_k, C_{k+1}) = S_\mathrm{stop} \cdot \mathbf{1}[C_k \ne C_{k+1}]
  + c_\mathrm{switch} \cdot |C_{k+1} \setminus C_k|.
\]

When $S_\mathrm{stop} = 0$ this reduces to the standard U-SSP.
When $c_\mathrm{switch} = 0$ (no individual switch cost) this is the TSIP (Section~\ref{sec:tsip}).
The interesting regime is $S_\mathrm{stop} > 0$ and $c_\mathrm{switch} > 0$.

\subsection{Collapse Condition}

\begin{thmbox}
\begin{theorem}[Setup-time collapse to JGP]
\label{thm:setup_collapse}
If the magazine setup time satisfies
\[
  S_\mathrm{stop} \;>\; n \cdot b \cdot c_\mathrm{switch},
\]
then SSP-$f_S$ collapses to JGP: the optimal solution uses exactly $\ZJGP$ groups.
\end{theorem}
\end{thmbox}

\begin{proofbox}
\begin{proof}
Consider two competing solutions: one with $K$ groups (configuration changes)
and one with $K+1$ groups.

The maximum possible total \emph{per-tool} switching cost across the entire
sequence of $n$ jobs and $b$-capacity magazine is bounded by $n \cdot b \cdot c_\mathrm{switch}$
(at most $b$ insertions per job, $n$ jobs total).

So any solution with $K$ groups costs at most $K \cdot S_\mathrm{stop} + n \cdot b \cdot c_\mathrm{switch}$.
Any solution with $K+1$ groups costs at least $(K+1) \cdot S_\mathrm{stop}$.

Under the condition $S_\mathrm{stop} > n \cdot b \cdot c_\mathrm{switch}$:
\[
  (K+1)S_\mathrm{stop} > K \cdot S_\mathrm{stop} + n \cdot b \cdot c_\mathrm{switch}
  \;\ge\; \text{cost of any } K\text{-group solution.}
\]
Therefore, no $(K+1)$-group solution can beat the best $K$-group solution.
Since this holds for any $K$, the optimal number of groups is minimised — exactly
the JGP objective. \qed
\end{proof}
\end{proofbox}

\begin{remark}[Interpretation]
When the setup time is large enough to dominate the per-tool cost of the
\emph{entire} sequence, it is always better to minimise stops (groups) than
to minimise individual tool insertions. The boundary $S_\mathrm{stop} = n \cdot b \cdot c_\mathrm{switch}$
is the \emph{collapse threshold}. Below this threshold, the SSP and JGP optima diverge.
\end{remark}

\subsection{Parameterised Collapse: The $\rho$-Frontier}

\begin{propbox}
\textbf{Status: Partially proved; approximation ratio analysis proposed.}
\end{propbox}

Define the ratio $\rho = S_\mathrm{stop} / c_\mathrm{switch}$.
As $\rho$ increases from 0 to $\infty$:
\begin{itemize}
\item $\rho = 0$: pure U-SSP; JGP and SSP optima can differ (gap analysis in
  \texttt{05\_jgp\_ssp\_gap\_analysis.tex}).
\item $\rho > nb$: complete collapse (Theorem~\ref{thm:setup_collapse}).
\item $0 < \rho < nb$: partial regime where the optimal number of groups
  is between $\ZJGP$ and some $K_\rho > \ZJGP$.
\end{itemize}

\noindent\textbf{Approximation ratio as a function of $\rho$.}
Define the JGP+GSP TD-SSP approximation ratio at parameter $\rho$:
\[
  R(\rho) = \sup_I \frac{(K^*-1)\rho + H(I)}{(K^*-1)\rho + \ZSSP(I)},
\]
where $H(I)$ is the JGP+GSP cost (U-SSP component only) and $\ZSSP(I)$ is the
U-SSP optimum. For the 6-job ring (\texttt{05\_jgp\_ssp\_gap\_analysis.tex}, Example~1),
$K^*=3$, $H=4$, $\ZSSP=3$:
\[
  R(\rho)_{\text{ring}} = \frac{2\rho + 4}{2\rho + 3}.
\]
At $\rho=0$: ratio $=4/3$. At $\rho=1$: ratio $=1$ (JGP+GSP is optimal).
For $\rho > 1 = \rho_c(\text{ring})$, the JGP+GSP solution is exactly optimal
for the TD objective.

\medskip
Two boundary values are established:
\begin{itemize}
\item $R(0) \leq 4/3$ for $b=3$ (conjectured tight; see
  \texttt{05\_jgp\_ssp\_gap\_analysis.tex}, \S5).
\item $\lim_{\rho \to \infty} R(\rho) = 1$ (JGP+GSP optimal in the TSIP regime).
\end{itemize}

\noindent\textbf{Practical implication.}
In CNC manufacturing, the machine stop time $S_\mathrm{stop}$ is typically
5--30 minutes (magazine swap, re-initialisation) while individual tool change
costs $c_\mathrm{switch}$ are 0.5--2 minutes, giving $\rho \approx 3$--60
\cite{Privault1995}.
For ring-like instances ($\rho_c=1$), all practical parameter values satisfy
$\rho \gg \rho_c$: JGP+GSP is exactly optimal for the TD objective.
The U-SSP gap (at $\rho=0$) is therefore largely irrelevant in practice for
such instances.

\medskip
\noindent\textbf{Bi-objective interpretation.}
The TD-SSP family parameterised by $\rho$ is the weighted-sum scalarisation of:
\begin{align*}
  \min\; & (K - 1) \quad &&\text{(configuration changes = JGP objective)},\\
  \min\; & Z_{\mathrm{switch}} \quad &&\text{(tool changes = U-SSP objective)}.
\end{align*}
As $\rho$ varies, the optimal solution traces the Pareto frontier.
JGP+GSP always achieves the JGP-efficient point $(K^*-1, H)$.
The U-SSP optimal sequence may achieve a different Pareto-efficient point
$(K > K^*-1, \ZSSP < H)$. For $\rho > \rho_c(I)$, the JGP point dominates.
The multi-objective structure is studied for related tool-switching variants
in \cite{Furrer2017}.

\noindent\textbf{Open problem.}
Characterise the exact threshold $\rho_c$ per instance, and the function $R(\rho)$
as a closed form. This is also related to the Lagrangian relaxation of
\texttt{06\_grouping\_selection.tex}, where $\lambda \approx \rho$ penalises extra
groups.

% ─────────────────────────────────────────────────────────────────────────────
\section{The TD-SSP and Lexicographic Collapse}
\label{sec:td_ssp}
% ─────────────────────────────────────────────────────────────────────────────

\subsection{Tool-Dependent SSP (TD-SSP)}

In the \emph{Tool-Dependent SSP} (TD-SSP), the cost of inserting tool $t$ is
$c_t > 0$ (not necessarily uniform). The transition cost is:
\[
  f_\mathrm{TD}(C_k, C_{k+1}) = \sum_{t \in C_{k+1} \setminus C_k} c_t.
\]

When $c_t = 1$ for all $t$, this reduces to the U-SSP. When $c_t$ are arbitrary,
the KTNS policy is no longer optimal (it assumes uniform costs). However, the
Tooling Problem (TP) for a fixed sequence remains polynomial via minimum-cost
network flow on an interval graph (totally unimodular constraint matrix)
\cite{Privault1995}.

\subsection{Lexicographic Collapse Under Heavy Penalties}

When a machine stop incurs a \emph{fixed} penalty $S_\mathrm{stop}$ in the
TD-SSP model, the total cost is:
\[
  Z_\mathrm{TD} = K \cdot S_\mathrm{stop} + \sum_{k=1}^{K-1} f_\mathrm{TD}(C_k, C_{k+1}).
\]

\begin{thmbox}
\begin{theorem}[Lexicographic JGP recovery for TD-SSP]
\label{thm:lexico}
If $S_\mathrm{stop} > n \cdot b \cdot \max_t c_t$, the TD-SSP objective is
\emph{lexicographically}: first minimise the number of groups $K$ (JGP objective),
then minimise the total tool-switching cost within those groups.
\end{theorem}
\end{thmbox}

\begin{proofbox}
\begin{proof}
The maximum possible total tool-switching cost (second objective) across the
entire sequence is bounded by $\sum_k \sum_{t \in C_{k+1} \setminus C_k} c_t
\le n \cdot b \cdot \max_t c_t$. Under the stated condition, this entire sum is
strictly less than $S_\mathrm{stop}$. Therefore, adding one extra machine stop
(increasing $K$ by 1) costs at least $S_\mathrm{stop}$, which strictly dominates
any possible savings in the per-tool cost. The primary objective is therefore to
minimise $K$ (JGP), and the secondary objective (per-tool cost) is resolved within
the $\ZJGP$-group structure. \qed
\end{proof}
\end{proofbox}

\begin{remark}[Practical relevance]
This theorem justifies the use of the JGP as a primary objective in settings
where machine downtime is very costly (e.g., aerospace manufacturing, high-precision
machining). In such settings, the JGP+GSP heuristic is \emph{exactly optimal} for
the modified objective, not merely an approximation.
\end{remark}

% ─────────────────────────────────────────────────────────────────────────────
\section{Analysis of Collapse Conditions}
\label{sec:analysis}
% ─────────────────────────────────────────────────────────────────────────────

\subsection{Which Conditions are Necessary?}

The conditions (C1)+(C2)+(C3) in Theorem~\ref{thm:sufficient} are
\emph{sufficient} for collapse, but not all are necessary.

\begin{thmbox}
\begin{proposition}[Necessity analysis]
\label{prop:necessity}
Of the three conditions (C1), (C2), (C3) in Theorem~\ref{thm:sufficient}:
\begin{enumerate}[label=(\roman*)]
\item \textup{(C3)} is necessary by definition of collapse.
\item \textup{(C1)} is \emph{not} necessary in general.
\item \textup{(C2)} is \emph{not} necessary in general.
\end{enumerate}
\end{proposition}
\end{thmbox}

\begin{proofbox}
\begin{proof}
\textbf{(C3) necessary.} If there exists a grouping with $K>\Kstar$ groups
whose SSP-$f$ cost is strictly lower than all $\Kstar$-group solutions, collapse
fails by definition. So (C3) is necessary.

\medskip
\textbf{(C1) not necessary.} Let $f(C,C')=\varepsilon>0$ for all pairs including
$C=C'$ (positive self-cost). Any solution with $K$ groups pays $(K-1)\varepsilon$
regardless of self-costs. Minimising $K$ still gives collapse, so (C1) fails
(self-cost positive) yet collapse holds.

\medskip
\textbf{(C2) not necessary.} Take $J=\{j_1,j_2,j_3\}$, $b=3$,
$T_1=\{t_1,t_2,t_3\}$, $T_2=\{t_3,t_4,t_5\}$, $T_3=\{t_1,t_4,t_6\}$.
Here $\Kstar=3$ (each pair requires $>3$ tools). Define:
$f(C_1,C_2)=0$, $f(C_2,C_3)=f(C_1,C_3)=1$, $f(C,C)=0$.
The best 3-group solution costs $f(C_1,C_2)+f(C_2,C_3)=0+1=1$.
No 4-group solution achieves cost $<1$ (it still must cross the boundary
between some pair with positive cost). Collapse holds despite $f(C_1,C_2)=0$
violating (C2). \qed
\end{proof}
\end{proofbox}

\begin{remark}[Minimal sufficient condition]
(C3) is the only necessary condition, but it is tautological (it states the
definition of collapse). The conditions (C1)+(C2)+(C3) form a clean and
\emph{verifiable} sufficient criterion. Finding a tight non-trivial
necessary-and-sufficient characterisation of all $f$ for which SSP-$f$
collapses to JGP remains open.
\end{remark}

\subsection{The Collapse Threshold $\rho_c(I)$}

For the setup-time model $f_S$ with ratio $\rho=S_\mathrm{stop}/c_\mathrm{switch}$,
the exact collapse threshold is computable per instance.

\begin{thmbox}
\begin{theorem}[Collapse threshold formula]
\label{thm:threshold}
Define the \emph{switch-savings function}:
\[
  \Delta(I,K) = \min_{\text{$\Kstar$-group solutions}} \mathrm{cost}_{c=1}(\cdot)
  - \min_{\text{$K$-group solutions}} \mathrm{cost}_{c=1}(\cdot),
  \quad K>\Kstar.
\]
The collapse threshold is $\rho_c(I)=\sup_{K>\Kstar}\Delta(I,K)/(K-\Kstar)$.
SSP-$f_S$ collapses to JGP if and only if $\rho>\rho_c(I)$.
\end{theorem}
\end{thmbox}

\begin{proofbox}
\begin{proof}
A $K$-group solution costs $(K-1)S_\mathrm{stop}+c_\mathrm{switch}\cdot z_K$.
It beats the $\Kstar$-group optimal iff
$c_\mathrm{switch}(z_{\Kstar}-z_K)>(K-\Kstar)S_\mathrm{stop}$,
i.e.\ $\rho<\Delta(I,K)/(K-\Kstar)$.
Collapse fails iff some $K$ satisfies this; collapse holds iff $\rho\ge\rho_c(I)$. \qed
\end{proof}
\end{proofbox}

\begin{remark}[The 6-job ring]
For the ring instance (Part~V, Example~1): $\Kstar=3$, best 3-group GSP cost
$z_3=4$, SSP optimum $z_4=3$ with 4 groups.
So $\rho_c(\text{ring})=\Delta(I,4)/(4-3)=(4-3)/1=1$.
For $\rho>1$: JGP is optimal (TSIP regime).
For $\rho\le1$: 4 configurations may be beneficial.
\end{remark}

% ─────────────────────────────────────────────────────────────────────────────
\section{Summary}
\label{sec:summary}
% ─────────────────────────────────────────────────────────────────────────────

\begin{center}
\begin{tabular}{llll}
\toprule
Variant & Cost $f$ & Collapse? & Status \\
\midrule
U-SSP & $|C_{k+1}\setminus C_k|$ & Not in general & Counterex.\ in Part~V \\
TSIP & $\mathbf{1}[C_k\ne C_{k+1}]$ & Yes, always & \textbf{Proved} \\
$f_S$, $S_\mathrm{stop}>nbc$ & Fixed stop + per-tool & Yes & \textbf{Proved} \\
TD-SSP, $S_\mathrm{stop}\gg c_t$ & Fixed stop + tool-dep. & Lexicographic & \textbf{Proved} \\
$f$ with C1+C2+C3 & General & Yes & \textbf{Proved (sufficient)} \\
\bottomrule
\end{tabular}
\end{center}

\bibliographystyle{abbrvnat}
\bibliography{references}
\end{document}
